% ============================================================================
% CAPÍTULO 7: ANTIMATERIA
% ============================================================================
\chapter{Antimateria}
\label{cap:antimateria}

\begin{center}
\textit{``El espejo de la materia''}
\end{center}

\vspace{0.5cm}

\begin{tcolorbox}[colback=white, colframe=kleinblue, boxrule=2pt]
\centering
\textit{``¿Por qué hay algo en lugar de nada?''}\\
--- Pregunta filosófica antigua\\[0.3cm]
\textit{``Porque $(7\pi)^{-7} \neq 0$.''}\\
--- Teoría Klein, 2026
\end{tcolorbox}


% ============================================================================
\section{El Espejo de la Materia}
% ============================================================================

Para cada partícula existe una ANTIPARTÍCULA:
\begin{itemize}
    \item Electrón ($e^-$) $\leftrightarrow$ Positrón ($e^+$)
    \item Protón ($p$) $\leftrightarrow$ Antiprotón ($\bar{p}$)
    \item Neutrón ($n$) $\leftrightarrow$ Antineutrón ($\bar{n}$)
\end{itemize}

Cuando materia y antimateria se encuentran: ¡ANIQUILACIÓN!
Toda la masa se convierte en energía pura ($E = mc^2$).

\textbf{EL GRAN MISTERIO:}

El Big Bang debió crear cantidades IGUALES de materia y antimateria.
¿Por qué el universo está hecho solo de materia?

La pequeña asimetría que permitió nuestra existencia:
\begin{equation}
\eta_B = \frac{n_B - n_{\bar{B}}}{n_\gamma} \approx 6 \times 10^{-10}
\end{equation}

Por cada mil millones de pares materia-antimateria que se aniquilaron,
sobró UNA partícula de materia. Esa es nuestra existencia.


% ============================================================================
\section{La Asimetría Bariónica: $\eta_B = (3/2) \times (7\pi)^{-7}$}
% ============================================================================

\begin{nivelmatematico}
\textbf{$\eta_B$ desde Klein}

\begin{equation}
\boxed{\eta_B = \frac{3}{2} \times (7\pi)^{-7}}
\end{equation}

donde:
\begin{itemize}
    \item $3/2$ = corrección cosmológica (3 dimensiones espaciales)
    \item $7$ = número de capas Klein (exponente)
    \item $(7\pi)^{-7}$ = supresión por atravesar 7 capas
\end{itemize}

\textbf{VERIFICACIÓN:}
\begin{align}
\text{Predicción:} & \quad \eta_B = 6.03 \times 10^{-10}\\
\text{Observado:} & \quad \eta_B = 6.12 \times 10^{-10}\\
\text{Error:} & \quad 1.5\%
\end{align}

\textbf{INTERPRETACIÓN:}

La asimetría es pequeña porque el proceso de bariogénesis
debe ``atravesar'' las 7 capas de la botella de Klein.
Cada capa suprime por factor $7\pi \approx 22$.
Total: $22^7 \approx 2.5 \times 10^9$
\end{nivelmatematico}


% ============================================================================
\section{Violación CP: $\varepsilon = (7\pi)^{-2}$}
% ============================================================================

\begin{nivelconceptual}
\textbf{Violación CP}

\begin{itemize}
    \item C = conjugación de carga (cambiar partícula por antipartícula)
    \item P = paridad (reflejar en un espejo)
    \item CP = ambas operaciones juntas
\end{itemize}

La física CASI respeta CP, pero no exactamente.
La violación CP es necesaria para explicar la asimetría materia-antimateria.
\end{nivelconceptual}

\begin{nivelmatematico}
\textbf{FÓRMULA KLEIN:}

\begin{equation}
\boxed{\varepsilon_{\text{CP}} = (7\pi)^{-2} \approx 0.00207}
\end{equation}

\textbf{VERIFICACIÓN:}
\begin{align}
\text{Predicción:} & \quad 0.00207\\
\text{Observado:} & \quad 0.00223\\
\text{Error:} & \quad 7.2\%
\end{align}
\end{nivelmatematico}


% ============================================================================
\section{La Oscilación Neutrón-Antineutrón}
% ============================================================================

Una predicción Klein aún no verificada:

\begin{equation}
\tau(n \to \bar{n}) \sim (7\pi)^{24} \times \tau_{\text{natural}} \sim 10^8 - 10^9 \text{ segundos}
\end{equation}

El neutrón puede ESPONTÁNEAMENTE convertirse en antineutrón.
Esto violaría la conservación del número bariónico.

El exponente $24 = \dim(\text{SU}(5))$ conecta con unificación de fuerzas.

Experimentos en ESS (European Spallation Source) buscan este efecto.


% ============================================================================
\section{Resumen del Capítulo}
% ============================================================================

\begin{tcolorbox}[colback=kleingray, colframe=kleinblue, title=\textbf{Resumen del Capítulo 7}]

\textbf{PREDICCIONES ANTIMATERIA KLEIN:}

\begin{enumerate}
    \item $\eta_B = (3/2) \times (7\pi)^{-7}$ \hfill [1.5\% error] --- Asimetría bariónica
    \item $\varepsilon_{\text{CP}} = (7\pi)^{-2}$ \hfill [7.2\% error] --- Violación CP
    \item $\tau(n \to \bar{n}) \sim (7\pi)^{24} \times \tau_{\text{nat}}$ \hfill [predicción] --- Oscilación $n$-$\bar{n}$
\end{enumerate}

\vspace{0.3cm}
\textbf{RESPUESTA AL MISTERIO:}

Existimos porque la bariogénesis atraviesa 7 capas Klein,
suprimiendo la simetría materia-antimateria por $(7\pi)^{-7}$.

\end{tcolorbox}


% ============================================================================
\section*{Código de Verificación}
% ============================================================================

\begin{lstlisting}[caption={Verificación numérica - Capítulo 7}]
import numpy as np
from numpy import pi

siete_pi = 7 * pi

# Asimetria barionica
eta_B_obs = 6.12e-10
eta_B_klein = (3/2) * siete_pi**(-7)
print(f"ASIMETRIA BARIONICA:")
print(f"  eta_B Klein = {eta_B_klein:.3e}")
print(f"  eta_B obs   = {eta_B_obs:.3e}")
print(f"  Error = {abs(eta_B_klein - eta_B_obs)/eta_B_obs*100:.1f}%")

# Violacion CP
epsilon_obs = 2.228e-3
epsilon_klein = siete_pi**(-2)
print(f"\nVIOLACION CP:")
print(f"  epsilon Klein = {epsilon_klein:.4f}")
print(f"  epsilon obs   = {epsilon_obs:.4f}")
print(f"  Error = {abs(epsilon_klein - epsilon_obs)/epsilon_obs*100:.1f}%")
\end{lstlisting}
